
\chapter*{Conclusiones}
\addcontentsline{toc}{chapter}{Conclusiones}

Los nodos implementados en el proyecto presentaron dificultades en la adaptación, dado que cada sensor entrega los datos según las condiciones del fabricante, por lo que se realizó el procesamiento de datos correspondiente para adecuarlos a los protocolos LoRaWAN.

El geolocalizador es uno de los sensores de mayor atención, en referencia a los utilizados en el nodo, dado que su sensibilidad permite determinar el lugar de la toma de datos. Si no hay geolocalización, el dato de interés se pierde, puesto que el nodo solo lo envía una o dos veces, de acuerdo con la clase de comunicación en la cual esté configurado.

El desempeño de este tipo de redes en la medición de parámetros provenientes de procesos, cuyo volumen de datos en el tiempo es bajo, queda demostrado con los levantamientos realizados durante las pruebas.

El nivel de pérdida de información es nulo, en parte por la pausa natural de los datos y, obviamente, por las técnicas aplicadas por el sistema. No se evidencian limitaciones respecto al tipo de variable ni al tipo de sensor empleado, en el contexto de monitoreo de parámetros ambientales.

La localización en interiores no cubrió las expectativas, puesto que LoRaWAN en clase A solo logró comunicaciones entre dos y tres pisos del edificio Sabio Caldas.

\enlargethispage{\baselineskip}

Los rangos de cobertura del sistema IoT LoRaWAN no alcanzan los valores indicados en las especificaciones del estándar, lo que puede atribuirse a que es una tecnología reciente. Las pruebas se realizaron con el sistema Multitech y con los equipos desarrollados en el GITUD.

La interoperabilidad entre los fabricantes (Multitech e ISMT) se comprobó mediante barridos entre la pasarela Multitech y el nodo ISMT, Mote II, el cual mostró una cobertura similar a la obtenida con el sistema Multitech.

El sistema LoRaWAN implementado en el grupo de investigación cumplió con las expectativas en cuanto a las comunicaciones y el funcionamiento de los nodos, logrando incluso una cobertura ligeramente superior a la del sistema Multitech, como se mostró en los mapas de cobertura.

Para futuros proyectos, se recomienda seleccionar sensores compatibles con los protocolos LoRaWAN, con el fin de minimizar el procesamiento de datos en el nodo. Asimismo, se recomienda emplear un geolocalizador de mayor sensibilidad para pruebas en interiores.

Con la tecnología LoRaWAN se pueden implementar sistemas IoT abiertos, dado que es compatible con \emph{software} y \emph{hardware} libres. Con la llegada de LoRaWAN clase C, es posible desarrollar sistemas inteligentes más robustos.

La ubicación de las pasarelas requiere de permisos institucionales, ya que, al tener antenas, deben instalarse en lugares elevados, con acceso a energía y conectividad. Para el caso de este proyecto, se consideró que la institución cuenta con varias sedes en Bogotá, lo que permitió realizar mediciones en dos puntos de la ciudad.

LoRaWAN clase A permite la configuración flexible de la red y, al ser un sistema abierto, puede adaptarse a diferentes tipos de aplicaciones y requisitos de la red, tanto libres como propietarios.

Los nodos LoRaWAN clase A solo pueden transmitir un número limitado de paquetes al gateway durante un periodo determinado. Esto puede representar una limitación para aplicaciones que requieren una alta tasa de transmisión de datos.

Si se requieren comunicaciones bidireccionales entre nodos y pasarelas, LoRaWAN ofrece las clases B y C, que permiten mayores tasas de datos que la clase A, aunque con una menor duración de la batería.
